\chapter{Гипотеза AM}
\noindent В главе рассматривается открытая гипотеза AM — равномерная по блокам нижняя оценка дискретного модуля путей через потерю радиуса корневой компоненты. Доказываются точная рекурсия этой потери, выделение первого разделяющего уровня и верхнее сравнение модуля с ценой конечного графового разреза. Недостающим остаётся именно равномерная нижняя оценка модуля; обсуждаются возможные геометрические условия для неё.

\section{Открытая оценка модуля путей}

\begin{openestimate}[Аннулярная оценка модуля путей (\(\mathsf{AM}\))]
\label{thm:annular-path-modulus}
Для голоморфно порождённых планарных Green-фильтраций с конечными
интерфейсными графами требуется доказать для всех достаточно больших
\(m\) с \(\Delta_m>0\) оценку
\eqref{eq:annular-modulus-target} с константой \(a_{x,r}>0\),
независящей от блока. В общей версии следует найти независимо
проверяемые геометрические гипотезы, не предполагающие локальную
связность предельного множества; естественные языки для их
формулировки — дискретный модуль, annular linear connectivity и
Loewner-type оценки.
\end{openestimate}
Гипотезы о Loewner-геометрии здесь не постулируются: необходимо
сформулировать независимо проверяемые условия, достаточные для
\eqref{eq:annular-modulus-target}, и затем проверить их на конкретной
фильтрации. Произвольные конечные графы для этого не подходят.

\section{Редукции к гипотезе AM: модуль путей}
\label{sec:am-open-reductions}

Открытая оценка $\mathsf{AM}$ требует равномерной по блоку нижней
границы
\[
 \operatorname{Mod}_{1,m}(\Gamma_m)\ge a_{c,r}\Delta_m,
 \qquad a_{c,r}>0,
\]
где $\Gamma_m$ состоит из всех путей через срез непосредственно до
первой потери корневой компоненты. Ни конечность интерфейсного графа,
ни одна лишь положительная дистанция между компонентами фиксированного
среза не дают константы, равномерной по глубине.

Пусть $c\in\partial\M$ и $r>0$. Для $n\ge1$ положим
\begin{equation}
 \label{eq:one-step-slice}
 A_n(c,r):=C_{n-1}(c,r)\cap p_n^{-1}(K),\qquad
 U_n(c,r):=C_{n-1}(c,r)\cap\{|p_n|>2\},
\end{equation}
где $\Sigma_n(c,r)$ определено в~\eqref{eq:green-separator}. Имеем
$C_n(c,r)=\operatorname{Comp}_c(A_n(c,r))$: всякое связное подмножество
$F_n(c,r)$, содержащее $c$, лежит в $C_{n-1}(c,r)$, а
$A_n(c,r)\subseteq F_n(c,r)$. Полуалгебраичность даёт конечное число
компонент у
$A_n$ и $U_n$; компоненты $A_n$ и замыкания компонент $U_n$ в
$C_{n-1}$ — компактные связные множества
\cite[теорема~2.4.5]{BCR}. По триангуляции полуалгебраических
множеств связное
$C_{n-1}(c,r)$ линейно связно; в частности, пути, используемые ниже,
существуют.

Положим
\begin{equation}
 \label{eq:block-newly-separated-set}
 \mathcal L_m(c,r):=\M\cap
 \bigl(C_{\lambda_m}(c,r)\setminus C_{\lambda_{m+1}}(c,r)\bigr),
 \qquad R_m:=\rho_{\lambda_m}^{\M}(c,r),\quad
 \Delta_m:=R_m-R_{m+1}.
\end{equation}
Из разбиения $\M\setminus C_n=(\M\setminus C_{n-1})\,\dot\cup\,
\bigl(\M\cap(C_{n-1}\setminus C_n)\bigr)$ следует точная рекурсия
\begin{equation}
 \label{eq:one-step-radius-recursion}
 \rho_n^{\M}(c,r)=
 \min\left\{\rho_{n-1}^{\M}(c,r),
 \dist\bigl(c,\M\cap(C_{n-1}(c,r)\setminus C_n(c,r))\bigr)\right\}.
\end{equation}
Итерация даёт
\begin{equation}
 \label{eq:block-radius-loss-exact}
 R_{m+1}=\min\{R_m,\dist(c,\mathcal L_m(c,r))\}.
\end{equation}
При $\Delta_m>0$ конечность числа компонент уровня
$O_{\lambda_{m+1}}\cap\overline B(c,r)$ показывает, что
$\mathcal L_m(c,r)$ есть конечное объединение пересечений $\M$ с
некорневыми компонентами этого уровня: всякая такая компонента,
встречающая $C_{\lambda_m}(c,r)$, целиком лежит в нём как связное
подмножество среза уровня $\lambda_m$. Поэтому $\mathcal L_m(c,r)$
компактно; выберем ближайшую к $c$ точку $y_m$.
Пусть $n_m$ — первый индекс в блоке, для которого
$y_m\notin C_{n_m}(c,r)$. Тогда
\[
 y_m\in C_{n_m-1}(c,r)\cap p_{n_m}^{-1}(K),\qquad
 y_m\notin C_{n_m}(c,r).
\]
Обозначим через $\Gamma_m$ семейство всех путей
$\gamma\in C^0([0,1],C_{n_m-1}(c,r))$ с
$\gamma(0)=c$, $\gamma(1)=y_m$; оно непусто по линейной связности
$C_{n_m-1}(c,r)$.
Поэтому всякий путь в $C_{n_m-1}(c,r)$ от $c$ до $y_m$ выходит из
$p_{n_m}^{-1}(K)$ и пересекает интерфейс $\Sigma_{n_m}(c,r)$.

Это разделение можно записать как конечный графовый разрез. Вершинами
графа служат компоненты $A_{n_m}$ и замыкания компонент $U_{n_m}$;
граф их пересечений связен, так как соответствующие компактные
связные множества покрывают связное $C_{n_m-1}$. Ребро между вершиной
$A\in\pi_0(A_{n_m})$ и замыканием $\overline U$ имеет контакт
\[
 P_e:=A\cap\overline U^{\,C_{n_m-1}}\subseteq\Sigma_{n_m}(c,r).
\]
Действительно, на $A$ выполнено $|p_{n_m}|\le2$, а на замыкании
$U$ — $|p_{n_m}|\ge2$; равенство следует из непустоты пересечения.
Положим
\[
 I_{n_m,m}:=\{\alpha:Q_{m,\alpha}\cap\Sigma_{n_m}(c,r)\ne\varnothing\}.
\]
Для $S\subseteq I_{n_m,m}$ удаляем ребро $e$ тогда и только тогда,
когда $P_e\subseteq\bigcup_{\alpha\in S}Q_{m,\alpha}$. Допустимым
называется такой $S$, который отделяет корневую вершину $C_{n_m}$ от
вершины $A^*$, содержащей $y_m$. Полный набор $I_{n_m,m}$ допустим,
поэтому минимум цены
$\sum_{\alpha\in S}\operatorname{diam}Q_{m,\alpha}$ достигается;
обозначим его через $\widehat{\mathfrak c}_m$.

Каждый допустимый разрез пересекает всякий путь
$\gamma\in\Gamma_m$. Действительно, прообразы вершинных множеств под
$\gamma$ дают конечное замкнутое покрытие $[0,1]$, граф пересечений
которого связен. Если $\gamma$ избегает выбранных ячеек,
то каждый контакт на этом графовом пути остаётся непокрытым и
соответствующее ребро не удаляется; тогда сохраняется путь от корня к
$A^*$, вопреки допустимости разреза. Следовательно, индикаторы
выбранных ячеек допустимы в определении дискретного модуля, и
минимум по $I_{n_m,m}$ совпадает с целевой ценой
$\widehat{\mathfrak c}_m$ из~\eqref{eq:modulus-below-cut}: ячейка,
покрывающая хотя бы одну точку контакта, сама встречает
$\Sigma_{n_m}(c,r)$.
\begin{equation}
 \operatorname{Mod}_{1,m}(\Gamma_m)
 \le\widehat{\mathfrak c}_m
 \le\operatorname{cap}_{n_m,m}
 \le\mathfrak c_m.
 \label{eq:am-cut-reduction}
\end{equation}
Второе неравенство следует из того, что целевой разрез отделяет только
одну $\M$-содержащую компоненту, а $\operatorname{cap}_{n_m,m}$
отделяет все такие компоненты; последнее — из определения полной
блочной цены.

Таким образом, строгая часть редукции даёт верхнюю оценку модуля через
цену разреза. Для замыкания аргумента нужна обратная, равномерная по
блоку нижняя оценка $\mathsf{AM}$. При фиксированном $n$ конечное число
компонент даёт положительную дистанцию от корневой компоненты до
остальных, но полученная константа зависит от $n$ и не контролирует
блоки при $m\to\infty$. Количественно достаточным языком остаются
дискретный модуль и равномерные annular/Loewner-оценки; они не
постулируются как доказанные свойства $\M$.

Для интерпретации этой недостающей равномерности можно ввести
минимальный радиальный обход
\[
 H_n(c,r;y):=\inf_{\substack{\gamma\in C^0([0,1],C_{n-1}(c,r))\\
                   \gamma(0)=c,\ \gamma(1)=y}}
                   \max_{t\in[0,1]}|\gamma(t)-c|.
\]
Если $y\notin C_n(c,r)$, каждый такой путь встречает $\Sigma_n$, откуда
$H_n(c,r;y)\ge\max\{|y-c|,\dist(c,\Sigma_n(c,r))\}$. Само это
топологическое пересечение не оценивает радиальный обход равномерно по
$n$; такая оценка была бы дополнительным путём к количественному
контролю потери.
